\begin{exercisechunk}
\item \textbf{Joint densities and conditional laws.}
\label[exercise]{ex:l5-continuous-em}
\exerciseprofile{3}{3}{3}{\exproof\quad\excalculation\quad\exconstruction}{%
Construct a conditional law from a common joint density and verify its
disintegration identity.}
Let \(Q_\theta\) be a family of probability laws on
\(\cX\times\cY\), where \(\cX\) and \(\cY\) are standard Borel spaces.
Fix sigma-finite reference measures \(P_X\) and \(P_Y\), put
\(P=P_X\otimes P_Y\), and assume that \(Q_\theta\ll P\). Write
\[
 q_\theta(x,y)=\frac{dQ_\theta}{dP}(x,y),
 \qquad
 m_\theta(y)=\int_\cX q_\theta(x,y)\,P_X(dx).
\]
Let \(Q_\theta^{(Y)}\) denote the marginal law of \(Y\).
\begin{enumerate}[(a),beginpenalty=10000]
\item Prove that \(Q_\theta^{(Y)}\ll P_Y\) and
\[
 \frac{dQ_\theta^{(Y)}}{dP_Y}(y)=m_\theta(y)
 \quad\text{for \(P_Y\)-almost every \(y\)}.
\]
\item For \(y\) such that \(m_\theta(y)>0\), define
\[
 r_\theta(x\mid y)=\frac{q_\theta(x,y)}{m_\theta(y)}.
\]
Show that this is a probability density with respect to \(P_X\). For
measurable \(A\subseteq\cX\), set
\[
 R_\theta(y,A)=\int_A r_\theta(x\mid y)\,P_X(dx).
\]
When \(m_\theta(y)=0\), define \(R_\theta(y,\cdot)\) to be any fixed
probability law on \(\cX\).
Verify that \(y\mapsto R_\theta(y,A)\) is measurable for each measurable
\(A\subseteq\cX\), and, for measurable \(B\subseteq\cY\), that
\[
 Q_\theta(X\in A,Y\in B)
 =\int_B R_\theta(y,A)\,Q_\theta^{(Y)}(dy).
\]
Conclude that \(R_\theta(y,\cdot)\) is a version of
\(Q_\theta(X\in\cdot\mid Y=y)\).
\end{enumerate}
\end{exercisechunk}
